% EPC flow charts
% Author: Fabian Schuh
\documentclass{article}
\usepackage{myflowchart}
\begin{document}

\begin{tikzpicture}

\begin{scope}[node distance=5mm and 5mm]

\node [ wideitem=2](a) at (1,1) {%
            \textbf{未完成方向}
            \nodepart[text width = 16.5cm]{two}
            \begin{enumerate}
\item a better frontpage arrangment in dbnote, apply it
\item event status section need minor correction.
\item Foolproof

\begin{enumerate}
	\item (upcoming page) export to html(blog) needs run str\_replace, because the greater than smaller than( google solution?)
\end{enumerate}

\item Bug
\begin{enumerate}
	\item when there is only grand final left, shouldn't show more match unfolding no?
	\item we are in between day 2 and 3 of ROG, day 2 completed, day 3 not entered, event not shown in schedule announced. : because day three has not been entered timely.( nothing can be done.... ignore)
	\item but want a link to go to the event page
	\item GO arrangment venture flow shows no git log?
	\item team logo not shown : there are two places of team logo, and the one matters need run collectstatic before one can commit the image, rather poorly designed 
	\item team page, eg liquid, has pass event should not shown epl s8: (this should be more easy now that countdown can be shown with timeuntil) because tournament and team views code still need to be adapted to the new regime, and need cleanup 
\end{enumerate}

\item Code Impovement :change path, incorporate os.path.join

\item Testing

\item wanted to do?  1. upcoming events, with no match schedule?  2. (CS) email steam? db model can be integrated into their system? big mission...
\end{enumerate}
            };


\node [above = of a, align = center] (title){\includegraphics[width=0.3\textwidth]{../../../figs/csgo_db_logo.png}\\ \LARGE GO事業部規劃};
            
\node [ wideitem=2, below = of a](unique_db) {%
            \textbf{資料庫賽程管理特色}
            \nodepart[text width = 16.5cm]{two}
            \begin{enumerate}
            	\item 能有個網頁可以馬上看到想看的結果，不用在上去Liquipedia或HLTV然後要點老半天找老半天，才找的到喜歡的隊伍的賽程。
            	\item 可選擇模組輸入或腳本輸入
            	\begin{enumerate}
            		\item 腳本輸入使用tts\_creation\_template.py檔。此步驟在隊伍的資料庫選擇上並沒有用上賽事單位的隊伍簡寫判斷，而是按照一排好的隊伍list，來做賽程輸入，若資料庫找不到隊伍，代表list中的隊伍名稱找不到對應的資料庫資料。這個狀況時，程式會允許使用者手動輸入正確的隊伍名稱。因此此方法適合比賽不多的賽事，不需要去跑格式判斷程式。編輯隊伍list的方法如下，從官網右鍵複製包含隊伍名稱及比賽時間的賽程表，貼上vim，用vim的快速組合鍵刪除所有不需要的東西，讓每一行剩下一個隊伍名稱，並且相連的兩個隊伍就是對戰的兩個隊伍。然後用我們note中整理的vim快速組合鍵將每個名稱加上"及",讓他變成一個python string list。編輯好後複製後就可以貼進tts\_creation\_template.py檔中。見我網站的db note page。
            		\item 模組輸入，網頁上複製包含該網頁所有隊伍及對戰時間的表格後，貼上一py檔，去跑test parsing.py，會自動分析格式取出隊伍名稱及對戰時間的程式。這對大型賽事較方便。如ESL pro League的巡迴賽，賽事有很多很多，就不適合手動輸入。目前在test\_scraping資料夾。
            	\end{enumerate}
           \end{enumerate}
            };


%\end{scope}
%\end{tikzpicture}
%
%%second page
%\begin{tikzpicture}
%\begin{scope}[node distance=5mm and 5mm]
\node [ wideitem=2, below =of unique_db](unique_editing) {%
            \textbf{網頁編輯特色}
            \nodepart[text width = 16.5cm]{two}
            \begin{enumerate}
            	\item upcoming post已整合成可用文書處理系統LYX編輯，不須處理raw html檔。
            	\item upcoming post目前是網路上少數可查找單一隊伍未來賽事的公開網頁，比HLTV、Liquipedia更方便。
           \end{enumerate}
            };

\end{scope}
\end{tikzpicture}
\newpage

%third page
\begin{tikzpicture}
\begin{scope}[node distance=5mm and 5mm]

	\node [ wideitem=2](gitlog) at (1,1) {%
            \textbf{last 10 git commits}
            \nodepart[text width = 16.5cm]{two}
            \footnotesize \verbatiminput{mysite.log}
            };

\end{scope}
\end{tikzpicture}

\end{document}
